% 04_lower.tex
\section{The lower side}
\label{sec:lower}

\subsection{The chain}
The proof is by contradiction. Fix a threshold $G$ and suppose $g<G$. Then, for any lower bound $F\le F_c$:
\begin{itemize}
\item[(L1)] By~\eqref{eq:branch}, $\sigma^2\Jww\le2A_w<\eta G\sigma\,\mathrm{den}$ and $N_\phi<\eta G\sigma\,\mathrm{den}$, with $\mathrm{den}\le\theta_G(b+W_p)$,
$b=\sigma^2/V$, $b+W_p=\sigma^2\Jvv\le1$.
\item[(L2)] The $w$-sandwich gives $L^2>\sigma/(\eta G\theta_G)$ and an upper bound on the window's $w$-information $u_0$: the window cannot be
narrower than about $\sqrt\sigma$ without its edges telling the observer $w$.
\item[(L3)] The curvature term with $\gamma_P$ and the AM--GM of its two parts force either a contradiction ($W_p\le b$) or
$b\le\beta W_p$, $V$ tiny, and $\sqrt{1-C}\,\omega<\Gamma_2\sqrt{W_p}$: the background's $v$-signal at the bent pendulum keeps the window's
$w$-extent short (with Lemma~\ref{lem:AY1P}, $L\le1$ is now a consequence).
\item[(L4)] Theorem~\ref{thm:AT2} with $\keff^2=xS+\sigma^2/(4V)$ and the $S$-bounds: $xS>F^2/(4G^2)-\tfrac14$, i.e. the ripple must carry a
share of the $w$-information.
\item[(L5)] Combining the three length relations with $q0_1^2\le4e^{-x}$ gives $e^{4x}<C(G,F)\,x^5/\sigma$, i.e. $x\le X_P$ with
$4X_P=\ls+5\ln X_P+K_P$. With (L4) and the Wiener bound $xS\le x+S1_P$: $\keff^2\le X_P+S1_P+\tfrac14$.
\end{itemize}
Then Theorem~\ref{thm:AT2} gives $g\ge F/(2\sqrt{X_P+S1_P+1/4})$.
In v1.1 every window lemma enters the chain in its $\theta_\tau$ form (App.~\ref{app:theta}): by Lemma~\ref{lem:BK1} a state with $g<G$ has
$\tau=k'^2V>80$, where the Gaussian averages differ from the period averages by a relative $\epsilon_\theta<10^{-13}$; the certificate perturbs every input
lemma by $\varepsilon_{\mathrm{in}}=10^{-6}$ in the bad direction (App.~\ref{app:cert}). The small constants ($\beta$, $\omega$, $\Gamma_2$, $\Gamma_3$, the range of $x$,
the case thresholds) are explicit functions of $(\sigma,G,\varepsilon,F,x)$; every supremum over $x$ is bounded cell by cell from
monotone factors (Appendix~\ref{app:cert}).

\subsection{The bootstrap for $F_c$}
The chain needs a lower bound $F\le F_c$. Lemma~\ref{lem:AY2P} supplies one, but its small quantities ($L^2$, $V$, $\sigma/s$) are bounded only
\emph{inside} the chain. We therefore read the argument as an upward bootstrap: from $F_c\ge F_k$ the chain gives $x\ge x_{\mathrm{low}}(F_k)$,
$L^2\le L^2_{\max}$, $s^2\ge\sigma/(\Gamma_{3P}x)$ and $x\le X_P$, and Lemma~\ref{lem:AY2P} then gives $F_c\ge F_{k+1}=(1-|\varepsilon|)F_{\mathrm{int}}(F_k)$.
Every step is an implication, so the iteration may only go up, from an \emph{assumed} start $F_0$. For periodic windows (with $x\ge1$) it closes
from $F_0=1.02G$: on the first cell ($10^{-24.5}\lesssim\sigma\le10^{-24}$) eight steps reach $F_{\mathrm{int}}\ge0.99955$.

\subsection{From points to a continuous range}
The constants of the chain depend on $\sigma$ in two ways: as a small factor (through $\eta G\sigma$, $\xi_P\propto\sigma$, $\sigma^2$, \dots) and through
$\ls$ (the a-priori range of $x$ and the root $X_P$); on paper~E's remainder only through the class bound $\bar\varepsilon$ of (P6),
via $\eta=1/(1-\bar\varepsilon)$ and the factor $1-\bar\varepsilon$.
Put $G(\ls):=\min(0.134,1.2/\sqrt{\ls})$ (independent of $\varepsilon$).

\begin{lemma}[monotonicity] \proofst
\label{lem:BH1}
Write $\hat\sigma$ for every occurrence of $\sigma$ as a small factor of the chain, and keep $\ls$ for its two occurrences as $\ln(1/\sigma)$ (the
a-priori range of $x$ and the root $X_P$). Every constant of the chain and of Lemma~\ref{lem:AY2P} is monotone in each of $(\hat\sigma,G,\eta,\ls,F)$,
in the directions of Table~\ref{tab:mono} (App.~\ref{app:cert}): $\theta_G$, $\beta$, $\nu$, $V_{\max}$, $\Gamma_2$, $\Gamma_{3P}$, $K_P$, $X_P$, $L^2_{\max}$,
$(\sigma/s)_{\max}$, $\delta_H$, $r_C$, $r_w$ do not decrease, and $\varphi_1$, $\omega$ do not increase, when $\hat\sigma$, $G$, $\eta$ or $\ls$ grows or $F$
decreases. $F$ enters only through $U_0=F^2/(4G^2)-\tfrac14$, which increases in $F$ and \emph{decreases in $G$}, and through
$x_{\mathrm{low}}=\max(1,U_0/S_i(1),U_0-S1_P)$, which does the same. Moreover $\eta$ enters the chain only through the upper bounds
$\sigma^2\Jww<\eta G\sigma\,\mathrm{den}$ and $N_\phi<\eta G\sigma\,\mathrm{den}$ of (L1), which remain true for every $\eta'\ge\eta$.
\end{lemma}
\begin{proof}
Table~\ref{tab:mono} writes each constant in $\hat\sigma$-normalised form (every $\sigma$-small factor appears with a positive power, and the
ratios that contain $\hat\sigma$ in numerator and denominator are reduced) and gives, row by row, the factor that carries each dependence.
The rows are composed from earlier rows in the order of the table, so the signs propagate: e.g.\ $\psi:=\eta G\theta_G$ increases in $(\hat\sigma,G,\eta)$
because $\theta_G=1/(1-\eta G\hat\sigma)$ does; $\varphi_1=1-e^{-1/4}-d_R$ decreases because $d_R$ is a sum of products of $\sqrt{\hat\sigma\psi}$, $\hat\sigma\psi$
and $S_\xi$; $\beta$ increases because $\varphi_1/(\varphi_1-2\psi^2)$ and $(2\psi+\hat\sigma)\psi/(\varphi_1-2\psi^2)$ do ($\varphi_1>8\psi^2$ by the $\Phi_C$
condition). The suprema over $x$ ($S_\gamma$, $S_\xi$, $S_\omega$, $S_q$, $S_N$) are suprema over $[x_{\mathrm{low}},x_e]$ whose right end $x_e$ grows with
$\ls$: they do not decrease in $\ls$ and do not increase in $x_{\mathrm{low}}$. The harmonic sums $\zeta_p$, $\Qb(x_{\mathrm{low}})$ decrease in $x_{\mathrm{low}}$.
\end{proof}

Hence, on a cell $\ls\in[\ell_a,\ell_b]$ and for all $\bar\varepsilon\le0.1$ at once, the chain may be evaluated at the worst corner: $\hat\sigma=e^{-\ell_a}$,
$\ls=\ell_b$ for the range, $G=G(\ell_a)\ge G(\ls)$ (the proof by contradiction from $g<G(\ls)$ also runs with any larger threshold), $\eta=1/0.9$, and the
bootstrap from $F_c\ge1.02G(\ls)\ge1.02G(\ell_b)$ (which every state with $g<G$ satisfies, Lemma~\ref{lem:BJ1}), with steps $F_{k+1}=0.9F_{\mathrm{int}}(F_k)$ (valid for every state with $|\varepsilon|\le0.1$ because
$F_c\ge(1-|\varepsilon|)A_1A_2A_3\ge0.9A_1A_2A_3$). The resulting $F_{\mathrm{int}}^*$ and $K^*$ hold for every state of the cell; inside the cell
$x\le X_P(\ls)$ with $4X_P=\ls+5\ln X_P+K^*$, and $\ls/(X_P(\ls)+c)$, $c=S1_P+\tfrac14$, is increasing in $\ls$ when
$d=4c-5-5c/X_P+5\ln X_P+K^*\ge0$ at $\ell_a$ (checked on every cell), so the cell's worst point is $\ell_a$ and no width is lost.

\begin{lemma}[tail beyond the certified range] \proofst
\label{lem:BH2}
For $\ls\ge\ell_{\mathrm{end}}=10^5\ln10$ (with $G=1.2/\sqrt{\ls}$ and the start $F_0/G=1.02$) all conditions of the chain hold, $F_{\mathrm{int}}(\ls)\to1$, and
$K_P(\ls)\le K_{\mathrm{end}}-\tfrac92\ln(\ls/\ell_{\mathrm{end}})$. Consequently
$\liminf_{\sigma\to0}g_{\min}^P\sqrt{\ls}\ge1-\bar\varepsilon$.
\end{lemma}
\begin{proof}
With $F/G$ fixed, Lemma~\ref{lem:BH1} makes every step at $\ls$ at least as good as at $\ell_{\mathrm{end}}$ ($\hat\sigma$ and $G$ are smaller; the suprema over $x$
do not depend on the right end of the range once $x_{\mathrm{low}}\ge8$, since all functions decrease beyond $x=8$). For $K_P$: $\Gamma_2\le(G/G_{\mathrm{end}})\Gamma_{2,\mathrm{end}}$
and $\Gamma_{3P}\le(G/G_{\mathrm{end}})\Gamma_{3P,\mathrm{end}}$ ($\theta_G$, $\beta$, $\zeta_2$ decrease, $\sigma/G$ decreases), $U_P\ge(G_{\mathrm{end}}/G)^2U_{P,\mathrm{end}}$ (from
$U_0=F^2/(4G^2)-\frac14$ with $F\ge F_{\mathrm{end}}$), and $\varphi_1$, $\omega$ do not decrease; so $K_P$ drops by $9\ln(G_{\mathrm{end}}/G)$. Every term of
$1-F_{\mathrm{int}}$ carries a factor $\sigma$, $\sqrt{\sigma X_P}$ or $L^2_{\max}\le Cx_{\mathrm{low}}e^{-x_{\mathrm{low}}}$ with $x_{\mathrm{low}}\ge x_{\mathrm{low,end}}\ls/\ell_{\mathrm{end}}$,
hence tends to $0$. Then $g\sqrt{\ls}\ge(1-\bar\varepsilon)F_{\mathrm{int}}\sqrt{\ls/(X_P+c)}/2$ with $X_P=\ls/4+O(\ln\ls)$, and the $G$-branch gives
$G\sqrt{\ls}=1.2$.
\end{proof}

\begin{remark}[the strip theorem read upward]
\label{rem:AY1}
The earlier strip certificate iterated $F$ downward from $F=1$, i.e. it used the fixed point ``$F_c\ge F\Rightarrow F_c\ge F_{\mathrm{int}}(F)\ge F$''. Read as a
proof it must start from an assumed $F_0$. For strips (no assumption $x\ge1$, so the first step has $x_{\mathrm{low}}<1$ and $L^2_{\max}=1$) the
upward bootstrap closes from $F_0=0.5$ at $\sigma=10^{-24}$ and from $F_0=0.2$ at $10^{-96}$ and below \caproofst{\code{bf\_ay\_bootstrap.py}}.
So the strip lower bound holds under the a-priori condition $F_c\ge\tfrac12$ (resp.\ $0.2$), with the same final constants:
$g\sqrt{\ls}/(1-|\varepsilon|)\ge0.921$ ($10^{-24}$), $0.977$ ($10^{-96}$), $0.998$ ($10^{-1024}$). A strip with $x\ge1$ is a periodic window and
needs no a-priori bound on $F_c$ (Lemma~\ref{lem:BJ1}, Theorem~\ref{thm:AY1P}). (The strip certificate is checked at these representative points only; the continuous
statement for strips with $x\ge1$ is Theorem~\ref{thm:AY1P}.)
\end{remark}

\subsection{The theorem}
\begin{theorem}[lower side] \caproofst{\code{bh\_cert\_cells.py}}
\label{thm:AY1P}
Let every state of the class satisfy (P0)--(P7) of Table~\ref{tab:assumptions}, with $|\varepsilon(\rho)|\le\bar\varepsilon\le0.1$. Then for every
$\ls\in[24\ln10,\,10^5\ln10]$ and every such state,
\begin{equation}
\begin{aligned}
g\ &\ge\ \min\Big(G(\ls),\ \frac{(1-\bar\varepsilon)\,F_{\mathrm{int}}^*}{2\sqrt{X_P+S1_P+1/4}}\Big),\\
\frac{g\sqrt{\ls}}{1-\bar\varepsilon}\ &\ge\ 1-c_1\,\frac{\ln\ls}{\ls},\qquad c_1=1.7,
\end{aligned}
\end{equation}
with $G(\ls)=\min(0.134,1.2/\sqrt{\ls})$, $4X_P=\ls+5\ln X_P+K^*$, $S1_P=1.41$, and $F_{\mathrm{int}}^*$, $K^*$ the certified constants of the cell
containing $\ls$ ($F_{\mathrm{int}}^*$ includes the factor $(1+\varepsilon_{\mathrm{in}})^{-1/2}$ of Lemma~\ref{lem:BL5}). The constants do not depend on $\bar\varepsilon$ (they are certified at $\bar\varepsilon=0.1$, Lemma~\ref{lem:BH1}). For $\ls>10^5\ln10$ the chain holds as
well and $\liminf g_{\min}^P\sqrt{\ls}\ge1-\bar\varepsilon$ (Lemma~\ref{lem:BH2}).
\end{theorem}

The certificate covers $[24\ln10,10^5\ln10]$ with $421$ cells of ratio $\ell_b/\ell_a=1.02$; every condition holds on every cell, and $d\ge8.9$.
The constant needed on a cell, $(1-L_{\mathrm{cell}})\ell_a/\ln\ell_a$, is at most $1.615$ (at $\ls\approx80.5$); with $c_1=1.7$ the smallest
absolute margin is $5.0\times10^{-5}$ (last cell; all with the input lemmas perturbed by $\varepsilon_{\mathrm{in}}=10^{-6}$). Table~\ref{tab:lower} shows representative cells.

\begin{table}
\caption{Theorem~\ref{thm:AY1P}: certified cell values $L_{\mathrm{cell}}\le g\sqrt{\ls}/(1-\bar\varepsilon)$, valid for every $\ls$ in the cell and every
$\bar\varepsilon\le0.1$ (representative cells of the $421$).}
\label{tab:lower}
\begin{ruledtabular}
\begin{tabular}{lccc}
cell containing & $X_P-\ell_a/4$ & $1-F_{\mathrm{int}}^*\le$ & $L_{\mathrm{cell}}$\\ \hline
$\sigma=10^{-24}$ & $+1.95$ & $4.5\times10^{-4}$ & 0.88999\\
$\sigma=10^{-48}$ & $+2.34$ & $5.7\times10^{-6}$ & 0.93466\\
$\sigma=10^{-96}$ & $+2.31$ & $9.1\times10^{-13}$ & 0.96593\\
$\sigma=10^{-256}$ & $+2.36$ & $6.5\times10^{-35}$ & 0.98650\\
$\sigma=10^{-1024}$ & $+2.50$ & $9.0\times10^{-142}$ & 0.99645\\
$\sigma=10^{-16384}$ & $+2.84$ & $4.9\times10^{-302}$ & 0.99975\\
$\sigma=10^{-100000}$ & $+3.06$ & $2.0\times10^{-302}$ & 0.99995\\
\end{tabular}
\end{ruledtabular}
\end{table}

The price of the comb's larger fundamental ($a_1$ up to $2$) is a constant $\ln256$ in $K_P$; the bootstrap more than compensates at
small $\sigma$ (it raises $x_{\mathrm{low}}$, which shrinks the harmonic sums). The bound rises to $1$ like $1-O(\ln\ls/\ls)$, which is why a fit over
a finite range of $\sigma$ returns a constant slightly different from $1$.
